\section{Implementation}

The current implementation is a C++20 interpreter built with CMake. The build
requires Flex and Bison. Bison generates the parser from
\code{src/platlang\_parser.y}, and Flex generates the lexer from
\code{src/platlang\_lexer.l}. The executable is named \code{platlang}.

At a high level, the interpreter follows a conventional pipeline. The lexer
recognizes keywords, identifiers, numbers, strings, comments, and punctuation.
The parser builds an abstract syntax tree with source locations. The
interpreter executes the AST using runtime value objects, table storage,
environments, function declarations, and localized diagnostic reporting.
Syntax errors stop execution before program evaluation. Runtime and semantic
errors are reported during interpretation.

Several implementation details support the design goals discussed above:

\begin{itemize}
  \item Top-level function declarations are registered before statements
        execute, allowing functions to call declarations that appear later in
        the file and allowing mutually recursive global functions.
  \item Variables are dynamically typed and stored in lexical environments.
        Assignment searches existing scopes, while \code{loat} creates a new
        binding in the current scope.
  \item Tables are mutable and use string or numeric keys. Dot access is a
        convenience for string-keyed fields.
  \item Diagnostics are represented by stable identifiers and rendered through
        localized templates for English and Limburgian.
  \item The command-line interface supports \code{--version}, \code{--ast},
        and diagnostic language selection through \code{--lang}.
\end{itemize}

The implementation is intentionally modest. It is valuable here because it is
inspectable and concrete, not because it introduces a new implementation
technique. That concreteness distinguishes \plat{} from a purely speculative
argument about localization: the design choices are embodied in a runnable
artifact with examples, tests, documentation, and diagnostics.

% TODO: Add an architecture figure in figures/ showing lexer, parser, AST,
% interpreter, environment, values, and diagnostics.
% TODO: Add a parser overview table after the grammar stabilizes.
% TODO: Summarize the runtime representation with references to source files.
% TODO: Add testing details, including unit tests and golden end-to-end tests,
% once the release branch is stable.
